% Solutions of Book 4, chapter 17 (the link layer).
\section*{Chapter \ref{ch:b2:link-layer} --- The Link Layer}

\begin{solution}{exo:b2:link-layer:1}
\code{0x4c} $= 0100\,1100$ has three 1s, so its bit is 1; \code{0x41} $= 0100\,0001$, \code{0x4d} $= 0100\,1101$ and
\code{0x50} $= 0101\,0000$ have two, four and two, so their bits are 0: \code{[1, 0, 0, 0]}, as printed. One flip in
each of two bytes is detected, in both bytes; three flips in one byte are detected (an odd number); two flips in one
byte are missed (\cref{prop:b2:link-layer:parity}).
\end{solution}

\begin{solution}{exo:b2:link-layer:2}
The loop adds the words but never folds \code{total}; the final \code{~total & 0xFFFF} then throws the
carries away instead of adding them back. For \code{LAMP}, $\mathtt{0x4c41} + \mathtt{0x4d50} = \mathtt{0x9991}$:
no carry, so both versions agree. For \code{ff f0 00 20}, $\mathtt{0xfff0} + \mathtt{0x0020} = \mathtt{0x10010}$; the
one's complement sum is $\mathtt{0x0010} + 1 = \mathtt{0x0011}$, whose complement is \code{ffee}, while dropping the
carry leaves \code{0x0010} and \code{ffef}. The bug shows only on data whose sum overflows, which is why a test on
short text missed it. The fix is the line \code{total = (total & 0xFFFF) + (total >> 16)} in the loop.
\end{solution}

\begin{solution}{exo:b2:link-layer:3}
\omcode{code/bachelor-2/17-link-layer/exercises/parity2d.py}{16}{33}{Correcting one flip with a parity grid}
\code{diff} has a 1 in each column whose parity is wrong. One wrong row and one wrong column: the data bit at their
crossing flipped, and XOR with \code{diff} puts it back. One wrong row and no wrong column: the row's \omterm{def:b2:link-layer:parity}{parity bit}
itself flipped, and the data is intact. No wrong row and one wrong column: a bit of the parity byte flipped. Anything
else needs two flips or more. The test flips each of the 32 data bits, each row bit and each column bit, and checks
that two flips in different rows give \code{None}.
\end{solution}

\begin{solution}{exo:b2:link-layer:4}
Append three zeros and divide, XORing $1011$ under each leading 1:
\[
11010011\,000 \to 01100011\,000 \to 00111011\,000 \to 00010111\,000 \to 00000001\,000 \to 00000000\,011 .
\]
The check bits are $011$ and the \omterm{def:b2:link-layer:frame}{frame} is $11010011\,011$. Dividing the \omterm{def:b2:link-layer:frame}{frame} repeats the same steps, whose last
three bits are now $011 + 011 = 000$: the remainder is 0. The test checks the value with \code{remainder}.
\end{solution}

\begin{solution}{exo:b2:link-layer:5}
\omoutput{code/bachelor-2/17-link-layer/outputs/exo_switch.txt}{The two tables after six \omterm{def:b2:link-layer:frame}{frames}.}
$A \to C$: both switches learn $A$ and flood ($s_1$ on ports 2 and 3, $s_2$ on 2 and 3). $C \to A$: $s_2$ learns $C$
on 2 and forwards on 1; $s_1$ learns $C$ on 3 and forwards on 1. $D \to B$: nobody knows $B$, so $s_2$ floods on 1
and 2, and $s_1$ on 1 and 2; both learn $D$. $B \to D$, $A \to B$ and $D \to C$ are forwarded, $A \to B$ by $s_1$
alone. Each switch sees every host behind the other switch on the port of the cable: \code{C:3 D:3} in $s_1$,
\code{A:1 B:1} in $s_2$. Two \omterm{def:b2:link-layer:frame}{frames} were flooded, by both switches.
\end{solution}

\begin{solution}{exo:b2:link-layer:6}
Stuffing sent \code{7e} as \code{7d 5e} and \code{7d} as \code{7d 5d}. The decoder drops each \code{7d} but leaves
the byte after it as it is, so \code{5e} and \code{5d} come out instead of \code{7e} and \code{7d}. It must XOR the
byte that follows an escape with \code{0x20}, as \code{unstuff} does. A second bug waits: a payload byte \code{7d}
would be dropped too, if it ever arrived unescaped from a faulty sender; \code{unstuff} treats it as an escape, which
at least keeps the rule in one place.
\end{solution}

\begin{solution}{exo:b2:link-layer:7}
\omcode{code/bachelor-2/17-link-layer/exercises/crc8.py}{5}{21}{CRC-8, bit by bit and by table}
With an 8-bit register and the message taken most significant bit first, each byte enters the register whole, so
\code{TABLE[b]} is the CRC of the single byte $b$ and the table step needs no shift: the state after a byte is the
table entry of the register XOR the byte. The test checks both against \code{remainder(m, 8n, 0x107)} on 40 random
messages.
\end{solution}

\begin{solution}{exo:b2:link-layer:8}
For a polynomial $P$, $P(1)$ is the sum of its coefficients, so it is the number of its terms modulo 2, and $(PQ)(1) =
P(1)Q(1)$, since evaluation respects products. Let $E$ be an error that passes the CRC: $g$ divides $E$, and $x + 1$
divides $g$, so $E = (x + 1)Q$ for some $Q$. Then $E(1) = (1 + 1)Q(1) = 0$: $E$ has an even number of terms, that is,
an even number of flipped bits. So every error with an odd number of flips is detected.
\end{solution}

\begin{solution}{exo:b2:link-layer:9}
\omcode{code/bachelor-2/17-link-layer/listings/backoff_sim.py}{5}{23}{\omterm{def:b2:link-layer:aloha}{Binary exponential backoff}}
\omoutput{code/bachelor-2/17-link-layer/outputs/backoff_sim.txt}{2000 runs per value of $k$.}
A \omterm{def:b2:link-layer:mac-protocol}{collision} costs one slot here, an idealisation of \omterm{def:b2:link-layer:aloha}{CSMA/CD}. The mean delay of a \omterm{def:b2:link-layer:frame}{frame} grows about linearly, close to
$2.4k$ slots: the stations spread out until the window is about as large as their number, and then go through
roughly one at a time. The time until the last \omterm{def:b2:link-layer:frame}{frame} grows faster, about 2.5 times per doubling of $k$: the unlucky
stations have collided many times and wait in windows much larger than needed. That unfairness is the known weakness
of exponential backoff.
\end{solution}

\begin{solution}{exo:b2:link-layer:10}
$\frac{d}{dG}\,G e^{-G} = (1 - G)e^{-G}$, positive for $G < 1$ and negative after: the maximum is at $G = 1$, and it is
$1/e \approx 0.368$. In pure \omterm{def:b2:link-layer:aloha}{ALOHA}, a \omterm{def:b2:link-layer:frame}{frame} sent at time $t$ survives if no other \omterm{def:b2:link-layer:frame}{frame} starts in $(t - 1, t + 1)$:
one starting just before overlaps its beginning, one starting just after its end. That interval has length 2 \omterm{def:b2:link-layer:frame}{frame}
times, so with Poisson attempts no other start falls in it with probability $e^{-2G}$, and the throughput is $G
e^{-2G}$. Its derivative is $(1 - 2G)e^{-2G}$, so the maximum is at $G = 1/2$ and equals $1/(2e) \approx 0.184$: half
of slotted \omterm{def:b2:link-layer:aloha}{ALOHA}'s, because slots halve the vulnerable interval.
\end{solution}

\begin{solution}{exo:b2:link-layer:11}
The proof of \cref{thm:b2:link-layer:crc-burst} shows that if $g$ divides $x C$ then it divides $C$; conversely, if
$g$ divides $C$ it divides $xC$. So $g$ divides $x^i B$ exactly when it divides $B$, whatever the position $i$.
\omoutput{code/bachelor-2/17-link-layer/outputs/burst_fraction.txt}{Every pattern of a \omterm{def:b2:link-layer:burst}{burst}, checked once.}
A \omterm{def:b2:link-layer:burst}{burst} of length 17 is $B$ of degree 16 with constant term 1. If $g$ divides it, $B = gQ$ with $Q$ of degree 0, so
$B = g$: one pattern in $2^{15}$. In \cref{out:b2:link-layer:bursts}, that one pattern sits at each of the $48 - 17 + 1
= 32$ positions, hence 32. For length 18, $Q$ has degree 1 and its constant term must be 1 (as $B(0) = g(0)Q(0)$), so
$Q = x + 1$: one pattern in $2^{16}$. In general, for \omterm{def:b2:link-layer:burst}{bursts} longer than $r + 1$ the fraction is $2^{-r}$.
\end{solution}

\begin{solution}{exo:b2:link-layer:12}
An error pattern on $N$ bits is a polynomial $E$ of degree less than $N$, and it passes when $E = gQ$. Each $Q$ of
degree less than $N - r$ gives such an $E$, and different $Q$ give different $E$; conversely the quotient of a
multiple of $g$ of degree less than $N$ has degree less than $N - r$. There are $2^{N-r}$ such $Q$, one of them 0, so
$2^{N-r} - 1$ nonzero patterns pass out of $2^N - 1$, about $2^{-r}$. If $x + 1$ divides $g$, every multiple of $g$
has an even number of terms (\cref{exo:b2:link-layer:8}), and the even patterns, $2^{N-1}$ of them (with 0), contain
all $2^{N-r}$ multiples: a fraction $2^{-(r-1)}$ of them pass, twice the average, while no odd pattern passes.
\end{solution}

\begin{solution}{pb:b2:link-layer:1}
\textbf{1.} \code{7e 7d 5d 7d 5e 00 7d 5e 7e}: each \code{7d} and \code{7e} becomes two bytes.

\textbf{2.} $2n$ bytes (plus the two flags), when every byte of the payload is \code{7e} or \code{7d}.

\textbf{3.} Two \omterm{def:b2:link-layer:frame}{frames} run together into one, of 12 bytes. The CRC does not catch it: it is accepted every time, not
with probability $2^{-16}$. Each \omterm{def:b2:link-layer:frame}{frame} $T_1$, $T_2$, check included, is a multiple of $g$, and with a register that
starts at zero and no final complement the merged \omterm{def:b2:link-layer:frame}{frame} is $T_1(x)\,x^{48} + T_2(x)$, a multiple of $g$ too, so its
check is 0. Only a length check (6 bytes expected) rejects it, or a CRC whose register starts at all ones, as real
\omterm{def:b2:link-layer:frame}{links} use: then the second \omterm{def:b2:link-layer:frame}{frame}'s check no longer matches what comes before it. (With one flag lost the \omterm{def:b2:link-layer:frame}{frames} stay
apart, since \code{stuffing.py} puts two flags between \omterm{def:b2:link-layer:frame}{frames}.)

\textbf{4.} The check protects the command, whatever the framing; the receiver removes the stuffing first, then
checks. The check bytes are stuffed like the payload, so a check that contains \code{7e} cannot end a \omterm{def:b2:link-layer:frame}{frame} early.

\textbf{5.} If the length byte is hit, the receiver reads the wrong number of bytes, takes the wrong place for the
next length, and every later \omterm{def:b2:link-layer:frame}{frame} is misread until something resynchronises it. A damaged flag costs only the \omterm{def:b2:link-layer:frame}{frames}
on both sides of it: the next flag is found again.

\textbf{6.} \code{LAMP} is \code{4c 41 4d 50} and its check is \code{2a 56}; none is \code{7e} or \code{7d}, so 6
bytes plus 2 flags: 8 bytes.

\textbf{7.} \code{4c 41 4d 50 2a 56}. The \omterm{def:b2:link-layer:frame}{frame} is $T = M x^{16} + R$, a multiple of $g$; the receiver's CRC is the
remainder of $T x^{16}$ by $g$, which is 0.

\textbf{8.} A \omterm{def:b2:link-layer:burst}{burst} often flips an even number of bits, which one \omterm{def:b2:link-layer:parity}{parity bit} never sees.

\textbf{9.} Every \omterm{def:b2:link-layer:burst}{burst} of 16 bits or fewer.

\textbf{10.} Every error with an odd number of flipped bits, since $g(1) = 1 + 1 + 1 + 1 = 0$, so $x + 1$ divides $g$.

\textbf{11.} A \omterm{def:b2:link-layer:burst}{burst} of 17 bits passes only if its pattern is $g$ itself (\cref{exo:b2:link-layer:11}): one pattern
per position, and there are $48 - 17 + 1 = 32$ positions.

\textbf{12.} $2^{24} \times (2^{32} - 1)/(2^{48} - 1) \approx 2^{24}/2^{16} = 256$. The simulation counted 256, all
with an even number of flips, as \cref{prop:b2:link-layer:crc-odd} demands. That it is exactly 256 is luck: the
standard deviation of such a count is about $\sqrt{256} = 16$.

\textbf{13.} Exactly one of the five sends: $5p(1 - p)^4$.

\textbf{14.} The derivative is $5(1 - p)^3(1 - 5p)$, zero at $p = 1/5$, where the value is $(4/5)^4 \approx 0.410$,
more than $1/e \approx 0.368$: with few stations the useful fraction $S_n = (1 - 1/n)^{n-1}$ at its best is larger,
and it decreases to $1/e$ as $n$ grows.

\textbf{15.} A radio cannot hear while it sends: its own signal is far stronger than any other at its antenna. And the
\omterm{def:b2:link-layer:mac-protocol}{collision} that matters happens at the receiver, which the sender may not hear at all.

\textbf{16.} The \omterm{def:b2:link-layer:hidden}{hidden terminal} problem: each finds the channel quiet, both send, and their \omterm{def:b2:link-layer:frame}{frames} collide at the
harbour, which receives neither; neither sender knows.

\textbf{17.} A sender first sends a short request; the harbour answers with a short ``clear to send'', heard by every
station near the harbour, the hidden one included, which then keeps quiet for the time announced. \omterm{def:b2:link-layer:mac-protocol}{Collisions} can now
hit only the short requests. The cost is two short \omterm{def:b2:link-layer:frame}{frames} and their delay for every data \omterm{def:b2:link-layer:frame}{frame}, which pays only for
long \omterm{def:b2:link-layer:frame}{frames}.

\textbf{18.} Noise and \omterm{def:b2:link-layer:mac-protocol}{collisions} lose far more \omterm{def:b2:link-layer:frame}{frames} on radio than on a cable, and the sender cannot detect a
\omterm{def:b2:link-layer:mac-protocol}{collision}; an acknowledgement on the \omterm{def:b2:link-layer:frame}{link} lets it resend at once, instead of leaving the loss to the much slower
layers above.

\textbf{19.} No: \omterm{def:b2:link-layer:burst}{bursts} of 17 to 20 bits each pass with probability about $2^{-15}$ or $2^{-16}$
(\cref{exo:b2:link-layer:11}). CRC-32 detects every \omterm{def:b2:link-layer:burst}{burst} of up to 32 bits.

\textbf{20.} Two more bytes per \omterm{def:b2:link-layer:frame}{frame}. With tables, the computation is one lookup per byte in both cases; the table
grows from 512 bytes to 1\,KiB.

\textbf{21.} $2/6 = 1/3$ with CRC-16; $4/8 = 1/2$ with CRC-32.

\textbf{22.} A \omterm{def:b2:link-layer:burst}{burst} of 20 bits passes with probability $2^{-16}$ (its pattern must be $gQ$ with $Q$ of degree 3 and
constant term 1: 4 patterns of $2^{18}$). Per million \omterm{def:b2:link-layer:frame}{frames}, 1000 have such a \omterm{def:b2:link-layer:burst}{burst}, and $1000 \times 2^{-16}
\approx 0.015$ pass: about one bad command in 65 million \omterm{def:b2:link-layer:frame}{frames}.

\textbf{23.} A table does in one lookup what takes eight shifts and tests; a hardware shift register with a few XOR
gates handles each bit as it arrives on the wire, at the speed of the \omterm{def:b2:link-layer:frame}{link}, with no memory at all.

\textbf{24.} Of 1\,114\,111 \omterm{def:b2:link-layer:burst}{bursts} of 1 to 16 bits, none passed CRC-16; of $2^{24} = 16\,777\,216$ random error
patterns, 256 passed, the predicted $2^{24} \cdot 2^{-16}$, all of them with an even number of flips.
\end{solution}
